\documentclass[12pt,a4paper]{article}
\usepackage[utf8]{inputenc}
\usepackage[T2A]{fontenc}
\usepackage[russian]{babel}
\usepackage{amsmath, amssymb}
\usepackage{geometry}
\geometry{top=2cm, bottom=2cm, left=2.5cm, right=2.5cm}
\usepackage{hyperref}
\hypersetup{colorlinks=true, urlcolor=blue}

\title{Метод GRA-обнулёнки для масштабирования дешёвых и доступных квантовых кубитов}
\author{Олег Битов}
\date{11 апреля 2026}

\begin{document}

\maketitle

\begin{abstract}
Современные квантовые компьютеры сталкиваются с фундаментальным ограничением: увеличение числа кубитов ведёт к экспоненциальному росту ошибок (декогеренция, шумы, тепловые флуктуации). Даже рекордный массив из 6100 кубитов требует сверхвысокого вакуума, температур, близких к абсолютному нулю, и сложных лазерных систем. В этой работе применяется многоуровневая GRA-обнулёнка к квантовым вычислениям. Вместо борьбы с ошибками за счёт всё более дорогой точности мы предлагаем иерархическую архитектуру, где дешёвые, но «пенящиеся» компоненты рекурсивно корректируются на нескольких уровнях. Показано, что такой подход ведёт к квадратичному уменьшению пены на каждом уровне, экспоненциальному подавлению ошибок и логарифмическим затратам ресурсов. Описан практический план создания миллиона кубитов стоимостью менее 1 миллиона долларов.
\end{abstract}

\section{Проблема: экспоненциальный рост ошибок при масштабировании}

Современные квантовые компьютеры сталкиваются с фундаментальным ограничением: при увеличении числа кубитов \textbf{пена} (ошибки декогеренции, шумы, тепловые флуктуации) растёт быстрее, чем возможности коррекции. Даже рекордный массив из 6100 кубитов требует сверхвысокого вакуума, оптических пинцетов и температуры, близкой к абсолютному нулю. Стоимость одного кубита остаётся огромной, а время когерентности ограничено (13 секунд – рекорд, но для вычислений нужно гораздо больше).

В терминах GRA-обнулёнки введём иерархию уровней:

\begin{itemize}
    \item \textbf{Уровень 0 (физический)}: отдельные атомы цезия, их квантовые состояния.
    \item \textbf{Уровень 1 (оптические пинцеты)}: лазерные ловушки, формирующие решётку.
    \item \textbf{Уровень 2 (система управления)}: лазеры, электроника, обратная связь.
    \item \textbf{Уровень 3 (квантовая коррекция ошибок)}: алгоритмы и резервные кубиты.
    \item \textbf{Уровень 4 (вычислительный)}: выполнение квантовых алгоритмов.
\end{itemize}

Пена \(\Phi^{(l)}\) на каждом уровне велика: ошибки позиционирования атомов, тепловые флуктуации, неидеальная изоляция, ограниченное время суперпозиции. Общий функционал:
\[
J_{\text{total}} = \sum_{l=0}^{4} \Lambda_l \Phi^{(l)} \gg J_{\text{крит}}.
\]

Традиционный подход пытается уменьшить каждую \(\Phi^{(l)}\) в отдельности (более мощные лазеры, лучший вакуум, более холодные температуры). Это экспоненциально дорого и упирается в физические пределы. GRA-обнулёнка предлагает иной путь: \textbf{создать новый уровень иерархии}, который рекурсивно снижает пену на всех нижележащих уровнях – подобно тому, как каменное орудие снизило пену физической уязвимости у австралопитеков.

\section{Принцип GRA-обнулёнки для квантовых кубитов}

Вместо того чтобы увеличивать количество кубитов прямым масштабированием (что ведёт к росту пены), предлагается \textbf{слоистая GRA-архитектура}, где кубиты организованы в иерархические кластеры с локальной обнулёнкой.

\subsection{Иерархия уровней кубитов}

Определим уровни \(l = 0, 1, \dots, L\):

\begin{itemize}
    \item \(l = 0\) (базовый кубит): физический атом в оптической ловушке.
    \item \(l = 1\) (логический кубит): группа из \(n_1\) базовых кубитов, связанных локальной GRA-коррекцией.
    \item \(l = 2\) (супер-логический кубит): группа из \(n_2\) логических кубитов.
    \item \(\dots\)
    \item \(l = L\) (вычислительный кубит): кубит, на котором выполняются операции пользователя.
\end{itemize}

Каждый уровень имеет свой \textbf{проектор} \(P_{G_l}\), который переводит состояние группы в ближайшее идеальное (без ошибок). Пена уровня \(l\):
\[
\Phi^{(l)} = \sum_{i,j} \bigl| \langle \Psi_i^{(l)} | P_{G_l} | \Psi_j^{(l)} \rangle \bigr|^2,
\]
где \(\Psi_i^{(l)}\) – состояния подсистем на уровне \(l\).

\subsection{Рекурсивная коррекция ошибок}

Вместо использования огромного числа резервных кубитов на одном уровне (как в поверхностных кодах), GRA-подход распределяет коррекцию по уровням. Каждый уровень исправляет ошибки только своего масштаба, передавая наверх только остаточную пену.

Это приводит к \textbf{логарифмическому} росту числа необходимых кубитов с глубиной иерархии:
\[
N_{\text{total}} = N_{\text{base}} \cdot L,
\]
где \(N_{\text{base}}\) – число базовых кубитов на уровень, а \(L\) – число уровней (логарифм от общего числа логических кубитов). Это колоссальная экономия по сравнению с поверхностными кодами, где требуется \(O(d^2)\) физических кубитов на один логический.

\subsection{Математическое обоснование}

Пусть вероятность ошибки на одном базовом кубите равна \(p\). Без коррекции, после \(m\) операций пена растёт как:
\[
\Phi^{(0)} \approx m p.
\]

При использовании GRA-иерархии из \(L\) уровней, где каждый уровень использует локальную коррекцию с эффективностью \(\eta\), остаточная пена после одного уровня:
\[
\Phi^{(l)} \approx \frac{1}{\eta} \left( \Phi^{(l-1)} \right)^2,
\]
то есть квадратичное уменьшение на каждом уровне. После \(L\) уровней:
\[
\Phi^{(L)} \approx \left( \frac{1}{\eta} \right)^{L} (m p)^{2^L}.
\]

Для достаточно больших \(L\) даже при умеренном \(p\) остаточная пена становится пренебрежимо малой. Это и есть \textbf{рекурсивное обнуление}.

\section{Практический метод создания дешёвых кубитов}

\subsection{Использование нетребовательных к вакууму и температуре атомов}

Вместо цезия (который требует сверхвысокого вакуума и наногельвиновых температур) можно использовать атомы рубидия или калия, которые могут быть захвачены в оптические решётки при менее экстремальных условиях. Однако у них выше пена (ошибки). GRA-иерархия компенсирует это: даже если \(\Phi^{(0)}\) на 2–3 порядка выше, чем у цезия, рекурсивная коррекция на следующих уровнях снижает её до требуемого уровня.

\subsection{Дешёвые лазерные пинцеты из обычных лазерных диодов}

Современные установки используют сложные твердотельные лазеры с высокой мощностью и стабильностью. GRA-подход позволяет заменить их массивом дешёвых лазерных диодов (как в принтерах), каждый из которых управляет небольшим кластером атомов. Пена от нестабильности каждого диода корригируется на уровнях 1 и 2.

Формула стоимости:
\[
C_{\text{total}} = N_{\text{кластеров}} \cdot C_{\text{диод}} + \text{накладные расходы на GRA-контроллеры}.
\]
Если кластер содержит 100 атомов, а лазерный диод стоит $10, то стоимость кубита падает до $0.1 плюс затраты на вакуумную камеру и электронику.

\subsection{Самообучение через GRA-оптимизацию}

Система непрерывно измеряет пену на всех уровнях и корректирует параметры управления (частоты лазеров, фазы, интенсивности) в реальном времени. Это аналог градиентного спуска из GRA-алгоритма. Со временем система находит оптимальные настройки, которые минимизируют \(J_{\text{total}}\), что позволяет использовать ещё более дешёвые компоненты.

\section{План развёртывания}

\begin{enumerate}
    \item \textbf{Фаза 1 (1–2 года)}: Создание прототипа с 10 кластерами по 100 атомов рубидия, управляемыми дешёвыми лазерными диодами. Демонстрация рекурсивной коррекции ошибок на 2 уровнях.
    \item \textbf{Фаза 2 (2–3 года)}: Масштабирование до 1000 кластеров (100 000 физических кубитов). Внедрение GRA-оптимизации в реальном времени. Достижение времени когерентности > 100 секунд (за счёт коррекции, а не изоляции).
    \item \textbf{Фаза 3 (3–5 лет)}: Создание модульной системы, где каждая плата содержит 10 000 кубитов. Соединение плат через оптические каналы. Общая стоимость системы – менее $1 млн за 1 миллион кубитов.
\end{enumerate}

\section{Выводы}

\begin{enumerate}
    \item \textbf{GRA-обнулёнка позволяет заменить дорогие высокоточные компоненты дешёвыми, но с высокой пеной}, за счёт рекурсивной коррекции ошибок на иерархических уровнях.
    \item \textbf{Математически} это выражается в квадратичном уменьшении пены на каждом уровне, что приводит к экспоненциальному подавлению ошибок при логарифмических затратах ресурсов.
    \item \textbf{Практически} это означает, что вместо одного массива из 6100 кубитов на сверхдорогой установке можно создать миллионы кубитов из доступных атомов и лазерных диодов, скомпонованных в GRA-иерархию.
    \item \textbf{Формула успеха}:
    \[
    \boxed{N_{\text{кубитов}} = \frac{C_{\text{бюджет}}}{C_{\text{диод}} \cdot \log \frac{1}{p_{\text{базовый}}}}}
    \]
    где \(p_{\text{базовый}}\) – допустимая пеня базового кубита. Чем выше \(p\) (дешевле компоненты), тем больше уровней иерархии нужно, но стоимость всё равно падает.
\end{enumerate}

Таким образом, GRA-обнулёнка открывает путь к \textbf{квантовому компьютеру для бедных} – точно так же, как она открывала путь к AGI для бедных, к освоению энергии галактики и к спасению человечества от самоуничтожения. Принцип един: не бороться с энтропией в лоб, а использовать её как ресурс, создавая новые уровни организации.

\end{document}